\section{Existence and absence of spurious right outliers}

Let $z_*$ be an isolated effective root of multiplicity $m$. Choose
$\delta_0>0$ so that
$[z_*-\delta_0,z_*+\delta_0]$ meets neither the bulk nor another root and
its endpoints are not roots. The count theorem gives exactly $m$
eigenvalues in this interval with probability tending to one. For every
$0<\eta<\delta_0$, the smaller interval
$[z_*-\eta,z_*+\eta]$ has the same limiting count $m$. On the intersection
of these two high-probability count events, all $m$ eigenvalues in the
larger interval lie in the smaller one. Hence all of them converge in
probability to $z_*$. This proves the existence assertion before it is used
below.

The spectrum also admits a deterministic high-probability bound. Indeed,
\begin{equation}
0\preceq M_d\preceq\frac{M}{n}A_dA_d^{\mathsf T}.
\end{equation}
Write $A_d=Z_d+\mathcal M_d$, where $Z_d$ has independent
$N(0,\lambda^{-1})$ entries and the columns of $\mathcal M_d$ are class
means. The Gaussian operator-norm inequality \cite{DavidsonSzarek2001}
\begin{equation}
\Pp\!\left\{
\|Z_d\|>
\lambda^{-1/2}(\sqrt n+\sqrt d+t)
\right\}
\le 2e^{-t^2/2}
\end{equation}
with $t=\sqrt d$ and $n/d\to\phi$ gives a deterministic $C_Z<\infty$ such
that $\Pp\{\|Z_d\|/\sqrt n\le C_Z\}\to1$. Moreover,
\begin{equation}
\frac{\|\mathcal M_d\|}{\sqrt n}
\le
\frac{\|\mathcal M_d\|_{\mathrm F}}{\sqrt n}
\le\max_b\|\mu_b^{(d)}\|
=
\max_b\sqrt{G_{C+b,C+b}}=:C_\mu.
\end{equation}
Thus, for
\begin{equation}
K_0=M(C_Z+C_\mu)^2,
\end{equation}
\begin{equation}
\Pp\{\|M_d\|\le K_0\}\longrightarrow1.
\label{S:norm-tight}
\end{equation}

Also, $S_G(z)=-z^{-1}+o(z^{-1})$ and therefore
\begin{equation}
F_G(z)\longrightarrow
\sum_{b=1}^kp_b\E[h_b(g)v_b(g)v_b(g)^{\mathsf T}]
\quad(z\to+\infty).
\end{equation}
Thus $F_G$ remains bounded while $zI_q$ diverges. There is a deterministic
$K_1$ beyond which $B_G(z)$ is positive definite, so all right effective
roots lie below $K_1$.

Fix $\epsilon>0$ and put $K=\max(K_0,K_1)+1$. We prove
\eqref{S:no-spurious} by its defining $\eta$-criterion. Let $\eta>0$ and
let
\begin{equation}
\mathcal Z_{\epsilon,K}
=
\{z\in[\lambda_++\epsilon,K]:\det B_G(z)=0\}.
\end{equation}
The set is finite. Choose pairwise disjoint neighborhoods of its points,
each of radius smaller than $\eta$. The compact remainder of
$[\lambda_++\epsilon,K]$ is a finite union of closed root-free intervals
after moving finitely many interval endpoints, if necessary, so that none
is a root. The count theorem gives no empirical eigenvalues in this
remainder with probability tending to one.

On the event $\|M_d\|\le K_0<K$, every empirical eigenvalue above
$\lambda_++\epsilon$ lies in $[\lambda_++\epsilon,K]$.  The preceding
compact-remainder argument therefore already accounts for the entire
relevant spectrum; no prior assignment of an empirical eigenvalue to a
limiting root is needed.  If $\mathcal Z_{\epsilon,K}\ne\varnothing$ then
\begin{equation}
\Pp\!\left\{
\sup_{\substack{x\in\sigma(M_d)\\x>\lambda_++\epsilon}}
\dist(x,\mathcal Z_{\epsilon,K})>\eta
\right\}\longrightarrow0,
\end{equation}
which is \eqref{S:no-spurious}. If
$\mathcal Z_{\epsilon,K}=\varnothing$, the compact-remainder argument
instead gives directly
\begin{equation}
\Pp\{\sigma(M_d)\cap(\lambda_++\epsilon,\infty)=\varnothing\}
\longrightarrow1,
\end{equation}
which is \eqref{S:no-spurious-empty}.


\section{Simple-root eigenvectors}

Let $z_d$ be an empirical eigenvalue converging to a simple root $z_*$ of
the original problem, let $u_d$ be a corresponding unit eigenvector, and
put
\begin{equation}
x_d=\frac{z_d}{\alpha_d},
\qquad
r_d=L_d^{\mathsf T}u_d.
\end{equation}
On the event $N_d\ge1$, \eqref{S:activeidentity} says that $x_d$ is an
eigenvalue of $\widehat M_d$ with the same eigenvector. The orthogonal
weighted-Wishart block has no eigenvalue in a fixed neighborhood of the
isolated active root with probability tending to one. On that event,
Lemma~\ref{S:exact-schur-lemma} gives the \emph{exact} identities
\begin{equation}
Q_{N_d}(x_d)r_d=0,
\qquad
r_d^{\mathsf T}Q_{N_d}'(x_d)r_d=1.
\label{S:active-exact-evec}
\end{equation}
Proposition~\ref{S:marked-theorem} and the random-index lemma give locally
uniform convergence of $Q_{N_d}$ and its derivative to the active
deterministic matrix $B_+$. Therefore
\begin{align}
\|B_+(x_d)r_d\|&=o_{\Pp}(1),\label{S:active-evec}\\
\left|r_d^{\mathsf T}B_+'(x_d)r_d-1\right|&=o_{\Pp}(1).
\label{S:active-evec-norm}
\end{align}

The exact deterministic scaling is
\begin{equation}
B_G(z)=\pi B_+\!\left(\frac z\pi\right),
\qquad
B_G'(z)=B_+'\!\left(\frac z\pi\right).
\label{S:B-scaling}
\end{equation}
It is not legitimate simply to replace the random $\alpha_d$ by $\pi$
inside the active equation. Using $z_d=\alpha_dx_d$, the required error is
decomposed exactly as
\begin{align}
B_G(z_d)r_d
={}&
\pi B_+(x_d)r_d
+(\alpha_d-\pi)x_dr_d \notag\\
&+\pi\left[
F_+(x_d)-F_+\!\left(\frac{\alpha_d}{\pi}x_d\right)
\right]r_d.
\label{S:evec-error-decomposition}
\end{align}
Here $\|r_d\|\le1$, the eigenvalues $x_d$ remain in a fixed compact right
set, and $F_+$ is locally Lipschitz there. Equations
\eqref{S:active-ratios}, \eqref{S:active-evec}, and
\eqref{S:evec-error-decomposition} prove \eqref{S:evec1}. The derivative
version follows from
\begin{equation}
B_G'(z_d)-B_+'(x_d)
=
B_+'\!\left(\frac{\alpha_d}{\pi}x_d\right)-B_+'(x_d),
\end{equation}
the local Lipschitz continuity of $B_+'$, and
\eqref{S:active-evec-norm}. This proves \eqref{S:evec2}.

By \eqref{S:multiplicity}, $\Ker B_G(z_*)$ is one-dimensional. On a compact
neighborhood of $z_*$, continuity gives $\sup_z\|B_G'(z)\|<\infty$; together
with \eqref{S:evec2}, this bounds $\|r_d\|$ away from zero with probability
tending to one. To make the final convergence precise, take an arbitrary
subsequence. Tightness in the compact unit ball gives a further subsequence
converging in distribution to a random vector $R$. Slutsky's theorem applied
to \eqref{S:evec1} and \eqref{S:evec2} gives, almost surely under the limiting
law,
\begin{equation}
B_G(z_*)R=0,
\qquad
R^{\mathsf T}B_G'(z_*)R=1.
\end{equation}
After the sign convention $\langle r_d,e_*\rangle\ge0$, these conditions
have the unique solution shown in \eqref{S:eveclimit}. Thus every
subsequential distributional limit is the same point mass. Convergence in
distribution to a constant is convergence in probability, proving
\eqref{S:eveclimit}.


\section{Multiple roots and spectral projections}
\label{S:multiple-section}

Let $z_*$ have multiplicity $m>1$. Choose a positively oriented contour
$\Gamma$ enclosing $z_*$ and no other effective root or bulk point. By
\eqref{S:pushforward} and \eqref{S:detscale}, $\Gamma/\pi$ encloses the
corresponding active root and is separated from the active bulk and all
other active roots.

Choose a sufficiently thin compact tubular neighborhood $U'$ of the
\emph{curve} $\Gamma/\pi$.  Because that curve has positive distance from
the active bulk and from every active root, one may choose an open annular
neighborhood $U\supset U'$ with $U'\Subset U$ whose closure meets neither
the active bulk nor a root.  The root enclosed by $\Gamma/\pi$ lies in the
interior of the contour, not in this annular neighborhood.  Hence $B_+^{-1}$
is holomorphic and bounded on $U$. For each deterministic active size,
Proposition~\ref{S:marked-theorem} and the exact compressed-resolvent
identity \eqref{S:exact-compressed-resolvent} give uniformly on $U'$,
\begin{equation}
L_d^{\mathsf T}(\widehat M_d-xI)^{-1}L_d
+B_+(x)^{-1}
\xrightarrow{\Pp}0.
\label{S:active-resolvent}
\end{equation}
Apply Lemma~\ref{S:random-index-lemma} to the event defined by the supremum
on $U'$. This transfers the uniform statement to the random active size
$N_d$. The exact scalar resolvent identity
\begin{equation}
L_d^{\mathsf T}(M_d-zI)^{-1}L_d
=
\alpha_d^{-1}
L_d^{\mathsf T}
\left(\widehat M_d-\frac z{\alpha_d}I\right)^{-1}
L_d
\label{S:resolvent-scaling}
\end{equation}
is valid whenever $N_d\ge1$. Since $\alpha_d\to\pi$, the random contour
$\Gamma/\alpha_d$ remains in the fixed neighborhood above. Moreover,
\begin{equation}
\alpha_d^{-1}B_+\!\left(\frac z{\alpha_d}\right)^{-1}
\longrightarrow
B_G(z)^{-1}
\end{equation}
uniformly on $\Gamma$, by \eqref{S:B-scaling}. Hence
\begin{equation}
\sup_{z\in\Gamma}
\left\|
L_d^{\mathsf T}(M_d-zI)^{-1}L_d+B_G(z)^{-1}
\right\|
\xrightarrow{\Pp}0.
\label{S:resolvent}
\end{equation}

Let
\begin{equation}
\Pi_d=-\frac1{2\pi i}\oint_\Gamma(M_d-zI)^{-1}\,dz
\end{equation}
be the full-space spectral projection. Integrating
\eqref{S:resolvent} gives
\begin{equation}
L_d^{\mathsf T}\Pi_dL_d
\longrightarrow
\frac1{2\pi i}\oint_\Gamma B_G(z)^{-1}\,dz.
\label{S:projection-contour}
\end{equation}
The kernel/complement expansion used in the multiplicity proof gives
\begin{equation}
B_G(z_*+t)^{-1}
=
\frac1tP_*J_*^{-1}P_*+O(1).
\end{equation}
Its residue is $P_*J_*^{-1}P_*$, proving \eqref{S:projection}. The count
theorem gives $\rank\Pi_d=m$ with probability tending to one. If $U_d$ is
any orthonormal basis of $\Ran\Pi_d$ and $R_d=L_d^{\mathsf T}U_d$, then
\begin{equation}
R_dR_d^{\mathsf T}
=L_d^{\mathsf T}\Pi_dL_d
\xrightarrow{\Pp}P_*J_*^{-1}P_*.
\end{equation}
The range of the limiting matrix is $K_*$, but an orthonormal basis of the
empirical eigenspace may rotate. No convergence of individual
eigenvectors is asserted.

Sections~\ref*{S:active-section}--\ref*{S:multiple-section} prove every item
of Theorem~\ref{S:main}: the active reduction and marked theorem give the
bulk and determinant, the random-index argument transfers the result to the
binomial active sample size, the exact rank identifies the zero modes, and
the final three sections give counts, exclusion of spurious outliers, and
the simple- and multiple-root eigenspace statements.

\section{ReLU-gate corollary}

\begin{corollary}
Assume
\begin{equation}
h_b(g)=a_b(g)\1_{\{\ell_b(g)>0\}},
\end{equation}
where $\ell_b$ is a nonconstant affine functional and
\begin{equation}
\tau\le a_b(g)\le M
\quad\text{on }\{\ell_b(g)>0\}.
\end{equation}
Then Theorem~\ref{S:main} applies.
\end{corollary}

\begin{proof}
A nonconstant affine functional of a standard Gaussian is a nondegenerate one-dimensional Gaussian, so its positive half-space has probability strictly between zero and one. The weights belong to $\{0\}\cup[\tau,M]$ almost surely.
\end{proof}

For the pure gate, fix a deterministic nonzero summary vector
$\theta\in\R^q$ and set $w_d=L_d\theta$ (equivalently, include the ReLU
normal among the fixed summary vectors used to construct $L_d$).  Then
\begin{equation}
w_d^{\mathsf T}Y_\ell
=
\theta^{\mathsf T}m_{b_\ell}
+\lambda^{-1/2}\theta^{\mathsf T}g_\ell,
\end{equation}
so the activation gate is exactly of the permitted form
$h_b(g)=\1_{\{\ell_b(g)>0\}}$ with
$\ell_b(g)=\theta^{\mathsf T}m_b+\lambda^{-1/2}\theta^{\mathsf T}g$.
In particular,
\begin{equation}
M_d=\frac1n\sum_{\ell=1}^n
\1_{\{w_d^{\mathsf T}Y_\ell>0\}}Y_\ell Y_\ell^{\mathsf T}.
\end{equation}
The exact neural interpretation is as follows. Fix a nonzero output coefficient
$a\in\R\setminus\{0\}$ and consider one ReLU unit
\begin{equation}
f_w(Y)=a\,\sigma(w^{\mathsf T}Y),
\qquad
\sigma(t)=t_+,
\end{equation}
with empirical squared loss
\begin{equation}
\mathcal L_n(w)
=
\frac1n\sum_{\ell=1}^n
\frac12\bigl(f_w(Y_\ell)-y_\ell\bigr)^2.
\end{equation}
At the fixed deterministic parameter $w=w_d$, every
$w_d^{\mathsf T}Y_\ell$ has a continuous distribution, so a finite
Gaussian sample avoids the kink almost surely. In a random neighborhood of
that fixed $w_d$, the activation pattern is constant and
\begin{equation}
\nabla_w f_w(Y_\ell)
=
a\,\1_{\{w^{\mathsf T}Y_\ell>0\}}Y_\ell,
\qquad
\nabla_w^2 f_w(Y_\ell)=0.
\end{equation}
Consequently, at that fixed non-kink parameter,
\begin{equation}
\nabla_w^2\mathcal L_n(w_d)
=
\frac{a^2}{n}\sum_{\ell=1}^n
\1_{\{w_d^{\mathsf T}Y_\ell>0\}}Y_\ell Y_\ell^{\mathsf T}.
\label{S:relu-hessian}
\end{equation}
The same matrix is the $w$--$w$ Jacobian-Gram or Gauss--Newton block.
Equation \eqref{S:relu-hessian} does not assert uniformity along a
data-dependent training trajectory, does not cover a parameter at a kink,
and does not describe cross-neuron, cross-layer, or generalized-Hessian
blocks.  A neuron normal with a component outside the fixed summary space is
also outside the present theorem unless that direction is first added to the
fixed-dimensional summary family.

